\documentclass{article}
\usepackage[utf8]{inputenc}
\usepackage[T1]{fontenc}
\usepackage{geometry}
\usepackage{graphicx}
\usepackage{float}
\usepackage{hyperref}
\usepackage[english,indonesian]{babel}

\title{Optimasi Pemrosesan Template LaTeX Menggunakan Pendekatan Berbasis AI}
\author{Academic Support Agent}
\date{\today}

\begin{document}
\maketitle

\begin{abstract}

\end{abstract}

\chapter{PENDAHULUAN}


\chapter{TINJAUAN PUSTAKA}


% Original content (commented out):
% %========== Preamble ==========
% 文档基本格式
\documentclass[a4paper,12pt,dvipsnames,UTF8]{article}

% 使用支持中文字体的宏包
\usepackage{ctex}

% 文本格式
\usepackage[normalem]{ulem}  % 文本强调
\usepackage{setspace} % 调整各种行间距
\interfootnotelinepenalty10000  % 规定脚注位置
\setlength{\parindent}{0cm} % 取消输入文本本身的缩进
\onehalfspacing % 1.5 倍行距

% 避免出现：①滞留前一页尾部的段落开头；②置于顶部的段落结尾。
\clubpenalty = 10000
\widowpenalty = 10000 
\displaywidowpenalty = 10000

% 数学公式
\usepackage{amsmath,amsfonts,amssymb} % 数学环境
\usepackage{array} % 打印矩阵
\usepackage{textcomp} % 其他数学字符

% 颜色库
\usepackage{xcolor}

% 图形图像
\usepackage{graphicx} % 用于插入图形
\usepackage{tikz} % 用于绘制简单形状
\usepackage{svg} % 用于插入.svg图片
\usepackage{float} % 图片排版工具
\usepackage[clockwise]{rotating} % 图片旋转工具

% 表格工具
\usepackage{tabularx} % 表格的设计
\usepackage{booktabs} % 经典三线表
\usepackage{multirow} % 更复杂的表格（例如“合并单元格”）
\newcolumntype{L}[1]{>{\RaggedRight\arraybackslash}p{#1}}
\newcolumntype{C}[1]{>{\Centering\arraybackslash}p{#1}}
\newcolumntype{R}[1]{>{\RaggedLeft\arraybackslash}p{#1}}

% 图表名称
\usepackage{caption}
\captionsetup{format=plain,justification=RaggedRight,singlelinecheck=false}
\usepackage{capt-of}
\usepackage[labelfont=bf]{caption}			
\captionsetup[table]{justification=centerlast}

% 页面格式
\usepackage[left= 3cm, right= 2cm, bottom= 2.5cm, top= 2.5cm]{geometry} % Side distances
\usepackage{pdflscape} % 页面旋转

% 图目录与表格目录
\usepackage{tocloft}
\renewcommand{\cftfigpresnum}{图}
\renewcommand{\cfttabpresnum}{表}

\renewcommand{\cftfigaftersnum}{:}
\renewcommand{\cfttabaftersnum}{:}

\setlength{\cftfignumwidth}{3cm}
\setlength{\cfttabnumwidth}{2,5cm}

\setlength{\cftfigindent}{0,5cm}
\setlength{\cfttabindent}{0,5cm}

% 文献引用
\usepackage{biblatex}
\addbibresource{Literature.bib}


%%%%%%%% 你的报告信息 %%%%%%%%
\newcommand{\titel}{主标题}
\newcommand{\untertitel}{副标题}
\newcommand{\autor}{作者A, 作者B, 作者C\\作者D, 作者E, ...}
%%%%%%%% 你的报告信息 %%%%%%%%


%========== 正文开始 ==========
\begin{document}

% 标题页
\input{Titlepage.tex}

% 摘要
\newpage % 另起一页
\pagenumbering{Roman} % 用罗马数字标记摘要与目录的页脚
\input{Abstract.tex}

% 目录页
\newpage % 另起一页
\tableofcontents % 打印目录（会由LaTex自动生成哦！）

% 图的目录（如果你需要的话）
\renewcommand\listfigurename{图目录}
\listoffigures
\addcontentsline{toc}{section}{图目录}

% 表格的目录（如果你需要的话）
\renewcommand\listtablename{表格目录}
\listoftables
\addcontentsline{toc}{section}{表格目录}

% 正文
\newpage % 另起一页
\pagenumbering{arabic} % 用阿拉伯数字标记正文页脚
\setcounter{page}{1} % 页码重置为1
\input{Outline.tex} %  设定文章提纲

% 参考文献（用.bib文件进行管理，例如本例的Literature.bib）
\newpage % 另起一页
\addcontentsline{toc}{section}{参考文献}
\printbibliography

% 附录（如果你需要的话）
\newpage % Creates a new page
\addcontentsline{toc}{section}{附录}
\input{Appendix.tex}

\end{document}
%========== 正文结束==========

\end{document}